\chapter{Charge sectors and the search for bound states}\label{chap:bound}

We now have all ingredients to answer the question of the introduction: which gauge-invariant particles does a
theory contain? After the gauge field has given a mass to some combinations of the bosons, the remaining
massless bosons are compact, and their charge sectors are the low-energy particles. Some of them are the original
fermions, others are \emph{composites}: gauge-invariant bound states of several fermions. The search is
implemented in \code{find_bound_states}.

\section{The idea}
A state with fermion numbers $N_L = (N_{L,1},\dots)$ and $N_R = (N_{R,1},\dots)$ carries gauge charge
\begin{equation}
  q_a = Q_{L,a}\cdot N_L + Q_{R,a}\cdot N_R = Q^V_a\cdot n + Q^A_a\cdot m ,
\end{equation}
with $n = N_L + N_R$ and $m = N_L - N_R$. On a ring the Gauss law $\oint\partial_xE = gq = 0$ allows only
\emph{neutral} sectors, $q_a = 0$ for all $a$. Neutral combinations of charged fermions are candidates for bound
states. Whether such a sector really contains a light particle is decided by its energy: on a ring of length $L$
a massless particle has energy $E\propto1/L$, so
\begin{equation}
  \Delta = \frac{EL}{2\pi} = h + \bar h
\end{equation}
stays finite as $L\to\infty$ for a massless sector, while a sector that only contains massive excitations does not
produce a $1/L$ level at all.

\section{The lattice computation}
On the ring every flavor has the zero modes $W_i = \sqrt\pi n_i$ and $P_{{\rm tot},i} = \sqrt\pi m_i$
(\cref{sec:ring}). The gauge field couples some combinations of them to the massive meson, so they are no longer
conserved; the conserved charges are the kernel directions of $MJ$, i.e.\ the anomaly-free global symmetries of
\cref{chap:anomalies}. The algorithm is:
\begin{enumerate}
  \item Build the ring Hamiltonian $M$ (\code{build_gauge_model} with \code{boundary="ring"}) and its conserved
    charges $A$ (\cref{sec:zero-modes}).
  \item For every neutral $(N_L,N_R)$ with $|N|\leq n_{\max}$: the configuration $\X_v$ with
    $W_i = \sqrt\pi n_i$, $P_{i,j} = \sqrt\pi m_i/N$ has conserved charges $\vect{c} = A^T\X_v$.
  \item Vectors with the same $\vect{c}$ are the same sector (they differ by excitations of the massive meson); keep
    the representative with the smallest $|N_L|^2 + |N_R|^2$.
  \item The lowest energy in the sector is $E = \frac12\vect{c}^TG\vect{c}$ \eqref{eq:constrained-minimum}: only
    the conserved charges are fixed, all other directions relax into the massive mode. The spin follows from the
    lattice momentum of the minimizing configuration, $s = -PL/2\pi$.
  \item Classify: $\Delta = |s| = \frac12$ is a single massless fermion; it is \emph{elementary} if the
    representative is one of the original fermions and a \emph{bound state} otherwise. All other sectors are
    multi-particle states.
\end{enumerate}
Because the zero-mode energies are exact on the lattice (\cref{chap:bosonization}), $\Delta$ is independent of $N$
and of $g>0$: the result is not an extrapolation.

\section{The analytic cross-check: the charge lattice}
The same answer follows from a short argument. The gauge fields give a mass to the combinations
$Q_{L,a}\cdot\varphi_L$ and $Q_{R,a}\cdot\varphi_R$ of the left- and right-moving bosons. What remains massless is
the part of $(N_L, N_R)$ orthogonal to the row spaces of $Q_L$ and $Q_R$:
\begin{equation}\label{eq:charge-lattice}
  h = \tfrac12\bigl|N_L^\perp\bigr|^2,\qquad \bar h = \tfrac12\bigl|N_R^\perp\bigr|^2,\qquad
  s = h - \bar h = \tfrac12\bigl(|N_L|^2 - |N_R|^2\bigr),
\end{equation}
for neutral $(N_L,N_R)$; the last equality uses neutrality. These vectors form the \emph{charge lattice} of the
low-energy theory (\code{analysis.charge_lattice}). \code{find_bound_states} computes both and reports whether
they agree (\code{SectorScan.agrees}).

\section{Worked example: the 3-4-5-0 model}
The 3-4-5-0 model~\cite{WangWen2019,Tong2022} has two left movers with charges $3,4$ and two right movers with
charges $5,0$ under one $U(1)$.
\begin{description}[leftmargin=1.5em,style=nextline]
  \item[Anomaly] $3^2 + 4^2 = 25 = 5^2 + 0^2$: consistent, although no fermion is vector-like.
  \item[Pairing] $(3,5)$ and $(4,0)$: $Q^V = (4,2)$, $Q^A = (-1,2)$; indeed $Q^V\cdot Q^A = 0$.
  \item[Meson] $\mu^2 = \frac{g^2}{2\pi}(9 + 16 + 25 + 0) = \frac{25g^2}{\pi}$, one massive and one massless boson.
  \item[Global symmetries] four fermion numbers, one condition \eqref{eq:mixed}: three conserved charges, one of
    them the gauge charge.
\end{description}
The search over $|N|\leq3$, $\Delta\leq2$ finds twelve sectors, in exact agreement with the charge lattice
(\cref{fig:sectors_3450}):
\begin{center}
\begin{tabular}{@{}lllll@{}}
\toprule
$N_L$ & $N_R$ & $(h,\bar h)$ & kind & content \\
\midrule
$(2,1)$ & $(-2,0)$ & $(\frac12,0)$ & bound state & $\psi_{3L}^2\psi_{4L}\psi_{5R}^{-2}$ \\
$(-2,-1)$ & $(2,0)$ & $(\frac12,0)$ & bound state & its antiparticle \\
$(0,0)$ & $(0,\pm1)$ & $(0,\frac12)$ & elementary & $\psi_{0R}^{\pm1}$ \\
four sectors & & $(\frac12,\frac12)$ & multi-particle & bound state $\times\ \psi_{0R}^{\pm1}$ \\
four sectors & & $(2,0)$, $(0,2)$ & multi-particle & two identical fermions \\
\bottomrule
\end{tabular}
\end{center}
The bound state is neutral, $2\cdot3 + 1\cdot4 - 2\cdot5 = 0$, and left-moving with $h = \frac12$: from
\eqref{eq:charge-lattice}, $N_L^\perp$ is the component of $(2,1)$ orthogonal to $(3,4)$, of length
$|4\cdot2 - 3\cdot1|/5 = 1$, and $N_R^\perp = 0$. It is a massless gauge-invariant composite fermion. Together with
the neutral right mover $\psi_{0R}$ it forms the massless low-energy content of the theory: one left- and one
right-moving fermion, central charges $(c_L, c_R) = (1,1)$ instead of $(2,2)$ before gauging. Counting particle and
antiparticle, the theory has two bound states.

\latsimfigure{sectors_3450}

\paragraph{Comparison with the literature.} The same four fermions carry a second anomaly-free $U(1)$,
$H = (5,0\,|\,3,-4)$, without mixed anomaly with $G = (3,4\,|\,5,0)$~\cite{Tong2022}; the library finds $H$ among the
surviving global symmetries of the gauged 3-4-5-0 model. Tong gauges $H$ and finds at low energies a single massless
Dirac fermion, the elementary left mover $\Psi = \psi_2$ and the gauge-invariant composite right mover
$\tilde\Psi = \psi_1^2(\tilde\psi_1^\dagger)^2\tilde\psi_2$, with $G$-charges $+4$ and $-4$~\cite{Tong2022}.
\code{find_bound_states(ChiralTheory([5, 0], [3, -4]))} reproduces exactly this: $\psi_2$ is elementary with
$(h,\bar h) = (\frac12,0)$, and the only bound states are $\psi_1^{\pm2}\tilde\psi_1^{\mp2}\tilde\psi_2^{\pm1}$ with
$(h,\bar h) = (0,\frac12)$ and $G$-charge $\mp4$. Gauging $G$ instead, as in this chapter, is the same theory up to
parity ($L\leftrightarrow R$) and charge conjugation of one fermion, and the sectors coincide with $h$ and $\bar h$
exchanged (\code{tests/test_literature.py}). Wang and Wen~\cite{WangWen2019} show that the 3-4-5-0 theory can be
realized as the low-energy limit of a local lattice model with an on-site $U(1)$ symmetry.

\paragraph{Volume dependence.} \Cref{fig:sector_volume_dependence} shows $EL/2\pi$ against the ring size. The
neutral massless sectors stay exactly at $h+\bar h$. Charged sectors are forbidden by the Gauss law on the ring; if
one computes them anyway, the seam of the ring acts as an external counter-charge that the massless matter screens
(\cref{sec:probes}), so $E$ stays finite and $EL$ grows: no massless particle.

\latsimfigure{sector_volume_dependence}

\section{Testing your own theory}
\begin{python}
from latsimpkg import ChiralTheory, find_bound_states
from latsimpkg.plotting import plot_sector_diagram, save_figure

theory = ChiralTheory(QL=[3, 4], QR=[5], g=0.3)   # rejects anomalous charges
print(theory.anomaly_matrix, theory.meson_masses, theory.global_symmetries().shape[1])

scan = find_bound_states(theory, N=40, nmax=3, Delta_max=2.0)
assert scan.agrees                                 # lattice = charge lattice
for s in scan.bound_states:
    print(s.content, s.h, s.hb)
save_figure(plot_sector_diagram(scan), "my_theory")
scan.save("my_theory.npz")                         # reload with SectorScan.load
\end{python}
In this example the right mover with charge $0$ is missing; the pairing adds it as a neutral spectator
(\code{spectator_R1}), and the bound states are those of the 3-4-5-0 model.

\paragraph{Limitations.}
\begin{itemize}
  \item The coupling must resolve the meson in double precision, $\mu L\gtrsim10^{-2}$ and $\mu a\lesssim10$
    (\cref{sec:numerics}); outside this window \code{find_bound_states} raises an error instead of returning sectors.
    The result itself does not depend on $g$: for 3-4-5-0 it is identical from $g = 10^{-4}$ to $10^2$.
  \item Only massless fermions with abelian gauge fields: the Hamiltonian must be quadratic. Massive fermions or
    four-fermion interactions add cosine terms to the bosonized Hamiltonian and require non-Gaussian methods.
  \item The search finds massless sectors. Massive bound states would be additional poles in correlation functions;
    in the Gaussian theories of this release the only massive particles are the free mesons.
  \item \code{nmax} limits which composites can be found: a bound state needs a representative with
    $|N_{L,i}|,|N_{R,i}|\leq n_{\max}$.
  \item The classification assumes a fermionic massless sector ($h\in k^2/2$), as for 3-4-5-0; for other
    theories inspect $(h,\bar h)$ directly.
\end{itemize}
